\section{Generalized harmonic numbers in every single-shift jet}
\label{asj:sec:harmonic}

Let $e_j(q)$ be the elementary symmetric polynomial of degree $j$ in
$1,1/2,\ldots,1/q$, with $e_0(q)=1$ and $e_j(q)=0$ for $j>q$.
For $q=0$ use the empty alphabet. Define
$H_q^{(j)}=\sum_{n=1}^q n^{-j}$ and $H_q=H_q^{(1)}$.
The finite product gives
\begin{align}\label{asj:eq:harmonic-egf}
 \frac{(s)_\ell}{\ell!}
 &=\frac{s}{\ell}\prod_{j=1}^{\ell-1}(1+s/j)
   =\frac{s}{\ell}\sum_{j=0}^{\ell-1}e_j(\ell-1)s^j,\notag\\
 \sum_{j\ge0}e_j(q)s^j
 &=\exp\left(\sum_{k\ge1}\frac{(-1)^{k+1}}kH_q^{(k)}s^k\right).
\end{align}
The second equality is understood as a power-series identity near zero;
its exponential simplifies to the finite product. In particular,
\[
 e_1(q)=H_q,\qquad
 e_2(q)=\tfrac12\{H_q^2-H_q^{(2)}\},\qquad
 e_3(q)=\tfrac16\{H_q^3-3H_qH_q^{(2)}+2H_q^{(3)}\}.
\]
These are the familiar harmonic/Stirling coordinates already present in
the manuscript's derivative tower. Their role here is to control
arithmetic spectral finite parts, not just parameter derivatives of a
single Hurwitz function.

\begin{theorem}[All single-shift Laurent jets]\label{asj:thm:harmonic}
For the multiplicative Hurwitz lattice, $|a|<\lambda_*$, and $r\ge1$,
\begin{align}\label{asj:eq:single-all}
 \calJ_r F_a
  ={}&\D^{(r)}(0)-a r!d_{r-1}\notag\\
   &+\sum_{\ell\ge2}\frac{(-a)^\ell}{\ell}
       \sum_{t=1}^{\min(r,\ell)}
         \frac{r!}{(r-t)!}\,e_{t-1}(\ell-1)
                                  \D^{(r-t)}(\ell).
\end{align}
The constant Laurent coefficient is
$\calJ_0F_a=\D(0)-a d_{-1}$. All spectral derivatives of $\D$
are finite sums of products of Hurwitz-zeta derivatives by Leibniz's
rule. Thus \eqref{asj:eq:single-all} uses only specified zeta jets,
generalized harmonic numbers, and finite Stieltjes polynomials.
\end{theorem}
\begin{proof}
The binomial expansion is
\[
 F_a(s)=\sum_{\ell\ge0}\frac{(-a)^\ell(s)_\ell}{\ell!}\D(s+\ell).
\]
The $\ell=0$ term contributes $\D^{(r)}(0)$. The resonant term $\ell=1$
is $-as\D(1+s)$ and contributes $-a r!d_{r-1}$, including $r=0$.
For $\ell\ge2$, substitute \eqref{asj:eq:harmonic-egf} and multiply by
the ordinary Taylor series of $\D(\ell+s)$. Theorem~\ref{asj:thm:alljets}
justifies the extraction and convergence.
\end{proof}

The first genuinely higher rows are worth displaying without compressed
symmetric-function notation:
\begin{align}\label{asj:eq:secondjet}
 \calJ_2F_a={}&\D''(0)-2a d_1
     +2\sum_{\ell\ge2}\frac{(-a)^\ell}{\ell}
           \{\D'(\ell)+H_{\ell-1}\D(\ell)\},\\
 \calJ_3F_a={}&\D'''(0)-6a d_2
     +3\sum_{\ell\ge2}\frac{(-a)^\ell}{\ell}
        \Big\{\D''(\ell)+2H_{\ell-1}\D'(\ell)\notag\\
        &\hspace{43mm}
        +(H_{\ell-1}^2-H_{\ell-1}^{(2)})\D(\ell)\Big\}.
 \label{asj:eq:thirdjet}
\end{align}
For two independent Hurwitz parameters, $d_1,d_2$ are exactly
\eqref{asj:eq:d1}--\eqref{asj:eq:d2}. For example,
\[
 \D'(\ell)=\zeta'(\ell,b)\zeta(\ell,c)
                    +\zeta(\ell,b)\zeta'(\ell,c),
\]
so every occurrence in \eqref{asj:eq:secondjet} is explicit. These formulas
are exact convergent identities; no assertion is made that the infinite
series always collapses to finitely many familiar constants.

\subsection{A finite algorithm and its shift continuation}
To compute a jet of order $r$, one first computes the finite Laurent
coefficients through $d_{r-1}$ using \eqref{asj:eq:drecipe}, then generates
the elementary symmetric coefficients through degree $r-1$, and finally
evaluates the normally convergent zeta-jet series. Only finitely many
algebraic operations are needed at each series index. The recurrence
\begin{equation}\label{asj:eq:Crecurrence}
 \ell C_\ell(\boldsymbol u)
     =\sum_{j=1}^{\ell}(-1)^j
         \left(\sum_i u_i a_i^j\right)C_{\ell-j}(\boldsymbol u)
\end{equation}
follows by differentiating \eqref{asj:eq:Cdef} with respect to $z$ and
is convenient for multiple factors. The verification code tests it
against the independent binomial product \eqref{asj:eq:Cpoch}.

The radius $|a|<\lambda_*$ is a convergence condition for this particular
Taylor series, not a singularity of the first-jet Gamma representation at
positive $a$. For larger shifts use the finite-head version of the transfer
theorem, or use \eqref{asj:eq:firstjet} and its derivatives. Confusing a
convergent expansion's radius with the full identity's domain would
unnecessarily discard most positive shifts.
